\documentclass[11pt,a4paper]{article}
\usepackage[provide=*,spanish]{babel}
\usepackage[utf8]{inputenc}
\usepackage[T1]{fontenc}
\usepackage{lmodern,amsmath,booktabs,tabularx,xcolor,listings,microtype,parskip,hyperref,titlesec}
\usepackage[margin=2.4cm]{geometry}
\definecolor{azulWCD}{RGB}{17,67,112}
\definecolor{tituloColor}{RGB}{0,55,110}
\definecolor{reglaColor}{RGB}{0,90,160}
\hypersetup{colorlinks=true,linkcolor=azulWCD,pdftitle={Densidades de los medios del WCD},pdfauthor={Jaime A. Betancourt}}
\lstset{language=C++,basicstyle=\ttfamily\small,breaklines=true,frame=single,backgroundcolor=\color{gray!6}}
\newcommand{\codigo}[1]{\texttt{#1}}
\titleformat{\section}{\normalfont\Large\bfseries\color{azulWCD}}
  {\thesection.}{0.6em}{}
\titleformat{\subsection}{\normalfont\large\bfseries}
  {\thesubsection.}{0.5em}{}

\begin{document}
\pagestyle{empty}

%------------------------------------------------
% PORTADA
%------------------------------------------------

\begin{titlepage}
\begin{center}

\vspace*{0.5cm}
\textcolor{reglaColor}{\rule{\linewidth}{1.2pt}}

\vspace{1.3cm}
{\Large\scshape Universidad Industrial de Santander}\\[0.2cm]
{\large Escuela de Física --- Facultad de Ciencias}\\[0.2cm]
{\large Doctorado en Física}

\vspace{2.2cm}
{\Huge\bfseries\color{tituloColor} Densidades de los\\[0.2cm]
medios del WCD}

\vspace{1.8cm}
{\large\textbf{Informe técnico}}

\vspace{0.6cm}
\begin{minipage}{0.82\linewidth}
\centering
\textit{``Implementación de la composición y densidad de los medios agua--NaCl en el WCD''}
\end{minipage}

\vfill

{\large\textbf{Autor}}\\[0.15cm]
Jaime A. Betancourt\\[0.1cm]
Doctorado en Física, Universidad Industrial de Santander

\vspace{1.2cm}
\textcolor{reglaColor}{\rule{0.5\linewidth}{0.6pt}}

\vspace{1cm}
11 de septiembre de 2026

\vspace{1cm}
\textcolor{reglaColor}{\rule{\linewidth}{1.2pt}}

\end{center}
\end{titlepage}

\pagestyle{plain}

\begin{abstract}
Se documentan las modificaciones propuestas para \codigo{Materials.cc} y
\codigo{SaltyWCD.cc} con el fin de representar, a $20\,^{\circ}\mathrm{C}$,
agua pura y soluciones acuosas con 2.5, 5.0 y 10.0\,\% de NaCl en fracción
másica final. Cada medio recibe su densidad indicada, la composición elemental
se calcula con la misma convención y la geometría selecciona el material
correspondiente. No se modifican la geometría, los modelos físicos de Geant4
ni las propiedades ópticas existentes.
\end{abstract}

\section{Convención de concentración}
La concentración se interpreta como fracción másica de NaCl respecto a la masa
total de solución:
\begin{equation}
r_{\mathrm{NaCl}}=\frac{m_{\mathrm{NaCl}}}{m_{\mathrm{agua}}+m_{\mathrm{NaCl}}},
\qquad
w_{\mathrm{agua}}=1-r_{\mathrm{NaCl}}.
\end{equation}
Por ejemplo, 2.5\,\% significa 2.5 g de NaCl y 97.5 g de agua por cada
100 g de solución. 
%Esta definición sustituye la interpretación anterior $m_{\mathrm{NaCl}}/m_{\mathrm{agua}}$, incompatible con la ecuación usada en la tesis.

\section{Densidades implementadas}
\begin{center}
\begin{tabular}{lcc}
\toprule
Medio & $w_{\mathrm{NaCl}}$ & Densidad (g cm$^{-3}$)\\
\midrule
Agua pura & 0.000 & 0.99820\\
Agua + 2.5\,\% NaCl & 0.025 & 1.01604\\
Agua + 5.0\,\% NaCl & 0.050 & 1.03401\\
Agua + 10.0\,\% NaCl & 0.100 & 1.07051\\
\bottomrule
\end{tabular}
\end{center}
Las densidades se asignan explícitamente y no se interpolan. Una concentración
no estudiada no recibe una densidad aproximada.

\section{Implementación en Materials.cc}
El agua pura se define como un material local, pues \codigo{G4\_WATER} usa la
densidad nominal $1.0\,\mathrm{g\,cm^{-3}}$:
\begin{lstlisting}
Water = new G4Material("Water",
                       WCDMediumDensity20C(0.0), 2);
Water->AddElement(elH, 2);
Water->AddElement(elO, 1);
\end{lstlisting}
La tabla óptica \codigo{waterPT1} se conserva.

Para las soluciones, \codigo{CreateSaltyWaterByFinalMassFraction} calcula:
\begin{align}
w_{\mathrm H}&=(1-r)\frac{2M_{\mathrm H}}{M_{\mathrm{H_2O}}},&
w_{\mathrm O}&=(1-r)\frac{M_{\mathrm O}}{M_{\mathrm{H_2O}}},\\
w_{\mathrm{Na}}&=r\frac{M_{\mathrm{Na}}}{M_{\mathrm{NaCl}}},&
w_{\mathrm{Cl}}&=r\frac{M_{\mathrm{Cl}}}{M_{\mathrm{NaCl}}}.
\end{align}
Las cuatro fracciones suman uno. La función
\codigo{WCDMediumDensity20C} acepta solo $r=0$, 0.025, 0.050 y 0.100; otro
valor genera \codigo{WCD-MAT-001}. También se admiten 2.5, 5 y 10 como
representación porcentual.

\section{Implementación en SaltyWCD.cc}
Este archivo no reconstruye la química. Lee
\codigo{detector.GetImpuritiesFraction()} y selecciona uno de los cuatro
materiales creados en \codigo{Materials.cc}. Se eliminó la aproximación al
medio más cercano: una concentración distinta genera
\codigo{WCD-MAT-002}. Así se evita simular silenciosamente, por ejemplo, una
entrada de 7\,\% con el medio de 5\,\%.

Durante la construcción se imprimen el nombre del medio, la fracción másica
normalizada y la densidad que Geant4 asignó. El volumen lógico del tanque usa
directamente ese \codigo{DetectorMedium}; por ello composición y densidad
cambian juntas al cambiar la configuración.

\section{Consecuencia física}
La densidad modifica la densidad numérica de H, O, Na y Cl por unidad de
volumen y, con ello, las longitudes libres medias y las probabilidades de
dispersión y captura. No debe aplicarse otra corrección de densidad en la
geometría: duplicaría o desacoplaría la propiedad almacenada en
\codigo{G4Material}.

\section{Verificación requerida}
Antes de producción se debe:
\begin{enumerate}
\item ejecutar \codigo{git diff --check};
\item compilar con MEIGA y Geant4 10.07.p04;
\item probar entradas 0, 0.025, 0.050 y 0.100;
\item repetir con 0, 2.5, 5 y 10;
\item verificar las cuatro densidades impresas;
\item comprobar que 0.07 sea rechazado;
\item comparar corridas cortas con la versión anterior.
\end{enumerate}
La revisión realizada es estática. La validación final requiere el repositorio
completo, la geometría, lista de física, datos nucleares y propiedades ópticas
de la investigación.

\section{Integración}
Los archivos se deben ubicar con nombres canónicos en
\codigo{src/Applications/G4WCDSimulator/}: \codigo{Materials.cc} y
\codigo{SaltyWCD.cc}. Las funciones nuevas son internas a
\codigo{Materials.cc}, por lo que no requieren declaraciones nuevas en
\codigo{Materials.h}. Si la cabecera conserva comentarios con la convención
anterior, deben actualizarse documentalmente.
\end{document}
